% =====================================================================
% Preamble: packages, styles, and custom commands
% Everything that configures the document lives here so main.tex stays clean.
% =====================================================================

% --- Encoding & fonts ------------------------------------------------
\usepackage[utf8]{inputenc}     % allow UTF-8 in source files
\usepackage[T1]{fontenc}        % proper output font encoding
\usepackage{lmodern}            % nicer-looking Latin Modern fonts
\usepackage{microtype}          % subtle typography improvements

% --- Page layout ------------------------------------------------------
\usepackage[margin=1in]{geometry}

% --- Math (used for things like phi X174) ----------------------------
\usepackage{amsmath, amssymb}

% --- Graphics and color ----------------------------------------------
\usepackage{graphicx}
\usepackage{xcolor}

% --- Hyperlinks (load near the end of packages, but before tcolorbox) -
\usepackage[colorlinks=true,
            linkcolor=blue!50!black,
            urlcolor=blue!50!black,
            citecolor=blue!50!black,
            pdftitle={Command Line for Biologists: A Practical Guide to the Terminal, Tools, and Reproducible Workflows},
            pdfauthor={Andrew Budge}]{hyperref}

% --- Code listings ---------------------------------------------------
\usepackage{listings}

% --- Colored callout boxes -------------------------------------------
\usepackage[most]{tcolorbox}

% --- Custom colors for code blocks -----------------------------------
\definecolor{bashbg}{HTML}{F7F7F2}
\definecolor{bashprompt}{HTML}{0066CC}
\definecolor{bashcomment}{HTML}{6A737D}
\definecolor{bashkeyword}{HTML}{B31D28}
\definecolor{bashrule}{HTML}{D0D7DE}

% --- Bash listing style ----------------------------------------------
% This styles every lstlisting block. The `literate` line recolors the
% `$` prompt to blue so terminal prompts stand out from the output.
\lstdefinestyle{bash}{
    backgroundcolor=\color{bashbg},
    basicstyle=\ttfamily\small,
    commentstyle=\color{bashcomment}\itshape,
    keywordstyle=\color{bashkeyword}\bfseries,
    breaklines=true,
    showstringspaces=false,
    frame=single,
    rulecolor=\color{bashrule},
    framesep=4pt,
    xleftmargin=8pt,
    xrightmargin=8pt,
    columns=fullflexible,
    keepspaces=true,
    aboveskip=1em,
    belowskip=1em,
    % Add your commands here as you introduce them in the book:
    morekeywords={pwd, ls, ll, cd, mkdir, touch, mv, cp, rm, rmdir,
                  echo, cat, less, head, tail, grep, find, chmod, chown},
    literate={\$\ }{{\textcolor{bashprompt}{\$}\ }}2,
}
\lstset{style=bash}

% --- Reflection callout box ------------------------------------------
% Use as:  \begin{reflection} ... \end{reflection}
\newtcolorbox{reflection}{
    colback=blue!5!white,
    colframe=blue!50!black,
    title=Reflection,
    fonttitle=\bfseries,
    boxrule=0.5pt,
    arc=2pt,
}

% --- Convenience commands --------------------------------------------
% \code{...} for inline code. Note: for code containing $, _, ~, \, #, %, &
% you must escape (e.g., \code{\$}) or use \verb|...| instead.
\newcommand{\code}[1]{\texttt{#1}}

% --- Spacing tweaks --------------------------------------------------
\setlength{\parskip}{0.5em}
\setlength{\parindent}{0pt}
